\documentclass[12pt]{article}
\usepackage{fullpage}
\usepackage{amsmath,amssymb,amsthm}

\newtheorem{lemma}{Lemma}
\newtheorem{proposition}{Proposition}
\newtheorem{assumption}{Assumption}
\newcommand{\R}{\mathbb R}
\newcommand{\E}{\mathbb E}
\newcommand{\1}{\mathbf 1}
\newcommand{\cE}{\mathcal E}
\newcommand{\Var}{\operatorname{Var}}

\begin{document}

\begin{center}
{\bf Invariant measures for the Dirichlet skew stochastic heat equation}
\end{center}

\section*{Problem statement and interpretation}

Let \(E=C_0([0,1])\) and let \(S_t\) denote the Dirichlet heat semigroup.
I use the ABLM regularized/mild interpretation specified in the problem.
The drift is the distributional derivative
\[
        \delta_0=-\frac12\frac{d}{dr}(-2H_+),\qquad
        H_+(r)=\1_{\{r>0\}} .
\]
Thus the natural candidate reversible measure is the bounded Gibbs
perturbation of the Brownian-bridge law.  Let \(\gamma\) be the law on
\(E\) of the standard Brownian bridge from \(0\) to \(0\), and set
\[
   \Lambda(v)=\int_0^1\1_{\{v(x)>0\}}\,dx,\qquad
   \pi(dv)=Z^{-1}e^{2\Lambda(v)}\gamma(dv).       \tag{1}
\]
The map \(\Lambda\) is Borel, since the evaluation map \((v,x)\mapsto
v(x)\) is Borel and Fubini applies.  Also \(0\le \Lambda\le1\), so
\(\pi\sim\gamma\), with density bounded above and below.  The value of
\(H_+(0)\) is irrelevant because a Brownian bridge spends zero Lebesgue
time at any fixed level:
\[
   \E_\gamma\int_0^1\1_{\{\beta(x)=0\}}\,dx
    =\int_0^1\mathbb P(\beta(x)=0)\,dx=0 .
\]

The argument below proves all parts of the standard reversible-measure
strategy which follow from the approximation assumed in the problem.  It
also isolates the single extra point needed for a complete uniqueness
proof.

\section*{What the approximation already gives}

Put
\[
   B_n(r)=\int_{-\infty}^rG_{1/n}(s)\,ds,\qquad
   \pi_n(dv)=Z_n^{-1}\exp\!\left(2\int_0^1B_n(v(x))\,dx\right)\gamma(dv).
\]
Then \(0\le B_n\le1\), \(B_n(r)\to H_+(r)\) for \(r\ne0\), and therefore
\[
       \|\pi_n-\pi\|_{\rm TV}\longrightarrow0 .          \tag{2}
\]
Indeed the normalized densities converge in \(L^1(\gamma)\) by dominated
convergence.

Let \(P_t^n\) be the Markov semigroup of the regularized equation with
drift \(G_{1/n}=B_n'\).  For smooth cylinder functions its generator is
\[
 L_nF(v)=\frac12{\rm Tr}\,D^2F(v)
       +\left\langle \frac12\Delta_Dv+G_{1/n}(v),DF(v)\right\rangle .
\]
The Gaussian integration-by-parts formula for \(\gamma\) gives
\[
   \int G\,L_nF\,d\pi_n
   =-\frac12\int\langle DF,DG\rangle_{L^2(0,1)}\,d\pi_n
   =\int F\,L_nG\,d\pi_n.                         \tag{3}
\]
Consequently \(P_t^n\) is reversible with invariant law \(\pi_n\).

\begin{proposition}\label{prop:limit}
Assume the approximation statement in the problem, in the following
standard consequence form: for each bounded continuous \(\Phi:E\to\R\),
each \(x\in E\), and each \(t>0\),
\[
        P_t^n\Phi(x)\longrightarrow P_t\Phi(x).      \tag{4}
\]
Then \(\pi\) is invariant for \(P_t\), \(P_t\) is symmetric on
\(L^2(\pi)\), and \(P_t\) has an \(L^2(\pi)\) spectral gap.
\end{proposition}

\begin{proof}
Let \(f,g\in C_b(E)\).  By dominated convergence with respect to the
fixed measure \(\pi\), by (2), and by invariance of \(\pi_n\),
\[
\begin{aligned}
 \int P_tf\,d\pi
 &=\lim_n\int P_t^nf\,d\pi
  =\lim_n\int P_t^nf\,d\pi_n
  =\lim_n\int f\,d\pi_n
  =\int f\,d\pi .
\end{aligned}
\]
Thus \(\pi\) is invariant.  The same argument, using (3), gives
\[
\begin{aligned}
 \int g\,P_tf\,d\pi
 &=\lim_n\int g\,P_t^nf\,d\pi_n
  =\lim_n\int f\,P_t^ng\,d\pi_n
  =\int f\,P_tg\,d\pi .
\end{aligned}
\]
Since \(C_b(E)\) is dense in \(L^2(\pi)\), and invariance makes \(P_t\) an
\(L^2(\pi)\)-contraction by Jensen's inequality, the symmetry extends to
all \(L^2(\pi)\).

It remains to pass a uniform Poincare inequality to the limit.  The
Brownian bridge satisfies
\[
     \Var_\gamma(F)\le C_0\int\|DF\|_{L^2(0,1)}^2\,d\gamma
\]
on cylinder functions.  Because the densities \(d\pi_n/d\gamma\) are
bounded above and below uniformly in \(n\), there is \(C<\infty\) such
that
\[
     \Var_{\pi_n}(F)\le C\,\cE_n(F,F),\qquad
     \cE_n(F,F)=\frac12\int\|DF\|_{L^2(0,1)}^2\,d\pi_n .       \tag{5}
\]
Hence the reversible semigroups \(P_t^n\) satisfy, with a constant
\(c>0\) independent of \(n\),
\[
     \Var_{\pi_n}(P_t^n f)
       \le e^{-2ct}\Var_{\pi_n}(f),\qquad f\in L^2(\pi_n).   \tag{6}
\]
For bounded continuous \(f\), (4) and (2) allow passage to the limit in
both sides of (6), yielding
\[
     \Var_{\pi}(P_t f)\le e^{-2ct}\Var_\pi(f).       \tag{7}
\]
Density and \(L^2(\pi)\)-contractivity extend (7) to every
\(f\in L^2(\pi)\).  This is the asserted spectral gap.
\end{proof}

\section*{The remaining uniqueness step}

The following elementary lemma is the point at which one needs a
pointwise smoothing or absolute continuity input for the actual ABLM
semigroup, not merely for a related \(L^2(\pi)\)-version.

\begin{lemma}\label{lem:criterion}
Assume, in addition to Proposition \ref{prop:limit}, that for some
\(\lambda>0\)
\[
     R_\lambda(x,A):=\int_0^\infty e^{-\lambda t}P_t(x,A)\,dt
        \ll \pi(A) \qquad\hbox{for every }x\in E .       \tag{8}
\]
Then \(\pi\) is the only invariant probability measure of \(P_t\).
\end{lemma}

\begin{proof}
Let \(\eta\) be invariant.  If \(\pi(A)=0\), then (8) and invariance give
\[
   \eta(A)=\lambda\int_E R_\lambda(x,A)\,\eta(dx)=0.
\]
Thus \(\eta=h\pi\).  For bounded \(\varphi\),
\[
 \int \varphi h\,d\pi=\int P_t\varphi\,h\,d\pi
   =\int \varphi\,P_t h\,d\pi ,
\]
where the last equality first holds for bounded \(h\) by symmetry and
then for \(h\in L^1(\pi)\) by truncation.  Hence \(P_t h=h\) in
\(L^1(\pi)\).

Fix \(a>0\), \(h_a=h\wedge a\).  Since \(r\mapsto r\wedge a\) is concave
and \(P_t\) is Markov,
\[
      h_a=(P_t h)\wedge a\ge P_t(h\wedge a)=P_t h_a .
\]
The two sides have the same \(\pi\)-integral, so equality holds
\(\pi\)-a.s.  Now \(h_a\in L^2(\pi)\), and the spectral gap (7) gives
\[
      \Var_\pi(h_a)=\Var_\pi(P_t h_a)
          \le e^{-2ct}\Var_\pi(h_a).
\]
For \(t>0\), this forces \(\Var_\pi(h_a)=0\).  Thus \(h\wedge a\) is
constant for every \(a>0\).  If \(0<a<b\), the constants for
\(h\wedge a\) and \(h\wedge b\), truncated again at \(a\), are equal;
letting \(a,b\) vary shows that \(h\) itself is constant.  Since
\(\int h\,d\pi=1\), \(h=1\), and \(\eta=\pi\).
\end{proof}

\section*{Relation with the Bounebache--Zambotti construction}

For the bounded variation potential \(f=-2H_+\), the
Bounebache--Zambotti Dirichlet form is precisely
\[
        \cE(F,G)=\frac12\int_E\langle DF,DG\rangle_{L^2(0,1)}\,d\pi .
\]
Their Proposition 2.5 supplies an \(L^2(\pi)\)-resolvent kernel
\(\rho_\lambda(x,dy)\) with \(\rho_\lambda(x,\cdot)\ll\pi\) for every
\(x\), and their associated Markov process solves the local-time form of
the skew heat equation for quasi-every initial point.  ABLM prove
well-posedness for \(b=\delta_0\) via regularized mild solutions; their
later weak/mild equivalence theorem gives a natural route for identifying
the local-time weak formulation with the ABLM formulation.

What is still missing for the present problem is a strict transfer of
that kernel statement to the particular semigroup \(P_t\) for every
deterministic \(u_0\in E\).  The approximation argument above avoids this
issue for invariance, symmetry and the spectral gap, but Lemma
\ref{lem:criterion} shows that uniqueness of arbitrary invariant
probabilities still requires the pointwise absolute continuity (8), or an
equivalent entrance-law statement.

\section*{Remaining open issues}

The proof is complete conditional on (8).  The unresolved point is:

\[
\boxed{\quad
   R_\lambda^P(x,\cdot)\ll\pi
   \quad\hbox{for every deterministic }x\in C_0([0,1]).
   \quad}
\]

Mathematically, this is needed only to rule out invariant probabilities
which might charge a \(\pi\)-null exceptional set.  This round established
that the actual ABLM semigroup is reversible with invariant law \(\pi\)
and has a spectral gap, using only the smooth approximations and the
total-variation convergence \(\pi_n\to\pi\).  The following possible
routes remain to close the last gap:

\begin{enumerate}
\item prove a uniform capacity or heat-kernel estimate for the
regularized semigroups \(P_t^n\) strong enough to pass
\(P_t^n(x,\cdot)\ll\pi_n\) to the limit;
\item prove that the ABLM solution is the strict version, for all starting points,
of the Bounebache--Zambotti Dirichlet-form process;
\item prove an entrance property saying that, even from a
Dirichlet-form exceptional point, the ABLM dynamics enters the
non-exceptional set immediately.
\end{enumerate}
Any one of these would imply (8), and then Lemma \ref{lem:criterion}
would prove that \(P_t\) has at most one invariant probability measure,
namely the measure (1).

\begin{thebibliography}{9}
\bibitem{ABLM}
S. Athreya, O. Butkovsky, K. L\^e and L. Mytnik,
\emph{Well-posedness of stochastic heat equation with distributional drift
and skew stochastic heat equation}, Comm. Pure Appl. Math. 77 (2024),
2708--2777.

\bibitem{ABLMweakmild}
S. Athreya, O. Butkovsky, K. L\^e and L. Mytnik,
\emph{Analytically weak and mild solutions to stochastic heat equation with
irregular drift}, Stoch. PDE: Anal. Comp. (2025).

\bibitem{BZ}
S. K. Bounebache and L. Zambotti,
\emph{A skew stochastic heat equation}, J. Theoret. Probab. 27 (2014),
168--201.
\end{thebibliography}

\end{document}
